\documentclass[border=8pt]{standalone}
\usepackage{fontspec}
\defaultfontfeatures{Ligatures=TeX}
\setmainfont{Palatino Linotype}
\setsansfont{Arial}
\usepackage{amsmath,amssymb}
% =====================================================================
% CAS PCBA Engineer — figures preamble (TikZ schematics)
% Little Henderson · Noetic Compass Press · 2026
%
% PREAMBLE-SAFE: this file is \input into the book preamble by the build.
% It must contain ONLY preamble code (package loads, color defs, \tikzset).
% No \begin{document}, no document-body material.
%
% Visual grammar (consistent across all 7 figures):
%   Domain colors as light fills + dark borders + black text:
%     P (Principled) = gold/amber, C (Cognitive) = blue,
%     E (Evaluative) = red,        O (Operational) = green.
%   Neutral / non-domain blocks = gray.
%   Rounded-rectangle blocks; clean downward signal-flow arrows (-{Latex}).
%   Designed to stay <= ~11cm wide (6x9 measure) and legible in grayscale
%   (rely on borders + labels, not color alone).
% =====================================================================

\usepackage{tikz}
\usetikzlibrary{arrows.meta,positioning,shapes.geometric,calc,fit,backgrounds}

% --- pgfplots for the vector charts (CJT curves, Bode, filters, decay) --
\usepackage{pgfplots}
\pgfplotsset{compat=1.18}
\usepgfplotslibrary{groupplots}            % stacked Bode panels share an x-axis
\pgfplotsset{
  % Shared house style for the analytic charts: light grid, sans labels,
  % measure-friendly default size, grayscale-legible line cycle.
  casaxis/.style={
    width=10.6cm, height=6.2cm,
    font=\sffamily\footnotesize,
    axis line style={draw=casNedge, line width=0.5pt},
    tick label style={font=\sffamily\scriptsize},
    label style={font=\sffamily\footnotesize},
    title style={font=\sffamily\footnotesize\bfseries},
    every axis grid/.style={draw=casNedge!25, line width=0.3pt},
    grid=both, tick align=outside, tickpos=left,
    legend cell align=left,
    legend style={font=\sffamily\scriptsize, draw=casNedge, fill=white,
      inner sep=2pt, row sep=0.5pt}
  }
}

% --- Domain palette (light fills, dark borders) ------------------------
\definecolor{casPfill}{RGB}{255,243,205}   % amber tint  (Principled)
\definecolor{casPedge}{RGB}{153,110,0}     % dark amber
\definecolor{casCfill}{RGB}{213,230,247}   % blue tint   (Cognitive)
\definecolor{casCedge}{RGB}{26,76,140}     % dark blue
\definecolor{casEfill}{RGB}{250,219,216}   % red tint    (Evaluative)
\definecolor{casEedge}{RGB}{150,40,30}     % dark red
\definecolor{casOfill}{RGB}{215,239,221}   % green tint  (Operational)
\definecolor{casOedge}{RGB}{30,110,55}     % dark green
\definecolor{casNfill}{RGB}{233,233,233}   % gray tint   (neutral)
\definecolor{casNedge}{RGB}{90,90,90}      % dark gray

% Accent for "go / pass" and "stop / refuse" paths (still labelled in text,
% so grayscale readers never depend on the color).
\definecolor{casGreen}{RGB}{30,110,55}
\definecolor{casRed}{RGB}{150,40,30}

% --- Reusable block + flow styles --------------------------------------
\tikzset{
  % Base block: rounded rectangle, black text, fixed measure-friendly width.
  casblock/.style={
    rectangle, rounded corners=2pt, draw, line width=0.6pt,
    text=black, align=center, inner sep=3pt,
    minimum height=8mm, font=\sffamily\footnotesize,
    fill=casNfill, draw=casNedge
  },
  % Domain-tinted blocks (inherit casblock, override fill+border).
  pblock/.style={casblock, fill=casPfill, draw=casPedge},
  cblock/.style={casblock, fill=casCfill, draw=casCedge},
  eblock/.style={casblock, fill=casEfill, draw=casEedge},
  oblock/.style={casblock, fill=casOfill, draw=casOedge},
  nblock/.style={casblock, fill=casNfill, draw=casNedge},
  % Wider container block for multi-line stage descriptions.
  widecasblock/.style={casblock, text width=58mm},
  widepblock/.style={pblock, text width=58mm},
  widecblock/.style={cblock, text width=58mm},
  wideeblock/.style={eblock, text width=58mm},
  wideoblock/.style={oblock, text width=58mm},
  widenblock/.style={nblock, text width=58mm},
  % Terminal nodes (signal entry / exit): small pill.
  casterm/.style={casblock, fill=white, draw=black, font=\sffamily\scriptsize\bfseries,
    minimum height=6mm, rounded corners=3pt},
  % Signal-flow arrow (downward by default).
  flow/.style={-{Latex[length=2mm,width=1.6mm]}, line width=0.6pt, draw=black},
  % Feedback / secondary path (dashed).
  fbflow/.style={-{Latex[length=2mm,width=1.6mm]}, line width=0.5pt, draw=casNedge,
    dashed},
  % Small italic edge label.
  flowlbl/.style={font=\sffamily\scriptsize\itshape, fill=white, inner sep=1pt},
  % Group / section frame drawn behind blocks (backgrounds layer).
  groupframe/.style={rounded corners=3pt, draw=casNedge, line width=0.5pt,
    dash pattern=on 1.5pt off 1.5pt, fill=casNfill, fill opacity=0.25},
  grouplbl/.style={font=\sffamily\scriptsize\bfseries, text=casNedge}
}

% Helper: render a CAS state literal in monospace with carets escaped, in
% math-free text mode. Usage: \casnote{^8^C|A}  ->  prints ^8^C|A safely.
\providecommand{\casnote}[1]{{\ttfamily\scriptsize #1}}

\begin{document}
\centering
\begin{tikzpicture}[x=1mm, y=1mm, font=\sffamily\footnotesize,
                    dnode/.style={rounded corners=1pt, minimum width=20mm,
                                  minimum height=8mm, draw, line width=0.6pt}]
  \node[dnode, fill=casPfill, draw=casPedge] (P) at (0,0)  {\textbf{P}};
  \node[dnode, fill=casCfill, draw=casCedge] (C) at (0,-14){\textbf{C}};
  \node[dnode, fill=casEfill, draw=casEedge] (E) at (0,-28){\textbf{E}};
  \node[dnode, fill=casOfill, draw=casOedge] (O) at (0,-42){\textbf{O}};
  % forward-biased diodes between domains
  \foreach \a/\b in {P/C, C/E, E/O} {
    \draw[-{Latex[length=1.8mm]}, casNedge, line width=0.6pt] (\a) -- (\b);
    \begin{scope}[shift={($(\a)!0.5!(\b)$)}]
      \draw[casNedge, line width=0.6pt] (-2.4,1.4) -- (2.4,1.4) -- (0,-1.4) -- cycle;
      \draw[casNedge, line width=0.6pt] (-2.4,-1.4) -- (2.4,-1.4);
    \end{scope}
  }
  \node[anchor=west, align=left, text=casNedge] at (14,-21)
    {forward bias:\\ conducts};
  % blocked reverse path
  \draw[casRed, line width=0.6pt, dash pattern=on 1.4pt off 1.4pt]
    (-13,-42) -- (-13,0);
  \draw[casRed, line width=0.7pt] (-15.6,-2.6) -- (-10.4,-7.8);
  \draw[casRed, line width=0.7pt] (-15.6,-7.8) -- (-10.4,-2.6);
  \node[anchor=east, align=right, text=casRed] at (-17,-21)
    {reverse bias:\\ blocked};
\end{tikzpicture}
\end{document}
